\newpage
\section{Deriving Optima}
The objective is minimum measured retrieval time at a required batch recall and a 32 GB
RAM limit. Only index and search parameters change. Document codes, query vectors and
reference answers remain identical. Each method chooses its own cluster count and local settings.

\subsection{Original method}
For a given layout, the reference answers show exactly how many clusters A must open.
Suppose a reference document belongs to the fifth cluster in the query's routing order.
A includes it when $P\ge5$. Record this rank for every reference document and query.

Let $s_{qi}^{(C)}$ be the cluster rank of reference result $i$ for query $q$. At recall
requirement $\rho$, the smallest qualifying probe count is
\[
 P_A(C,\rho)=\min\left\{P:
 \frac{1}{n_qK}\sum_{q=1}^{n_q}\sum_{i=1}^{K}
 \mathbf 1[s_{qi}^{(C)}\le P]\ge\rho\right\}.
\]
The indicator contributes one when the document's cluster is opened. With 200 queries
and 100 reference IDs each, 99\% recall requires at least 19,800 recovered IDs across the
batch. These ranks are calculated before timing the searches.

The experiment uses one global cluster and learned layouts with 256, 4096 and 8192 clusters.
For each layout it computes probe cutoffs for 80\%, 90\%, 95\% and 99\% recall, and also opens
all clusters. Query time and selected-row counts are then measured. These cutoffs describe
this query batch; they do not guarantee the same recall on future queries.

A derivative explains the competing costs. Suppose clusters are balanced, routing costs
$\alpha C$, each row costs approximately $\beta$, and $P$ stays constant. Then
\[
 T_A(C)\approx\alpha C+\beta NP/C,\qquad
 \frac{dT_A}{dC}=\alpha-\frac{\beta NP}{C^2},\qquad
 C^*=\sqrt{\frac{\beta NP}{\alpha}}.
\]
This balances routing more clusters against scanning fewer rows. Under a recall requirement,
$P_A(C,\rho)$ can change with $C$. The measured comparison uses that probe count and actual
cluster sizes; the fixed-$P$ derivative is not a complete solution.

\newpage
\subsection{Bitplanes}
Consider a mask covering 6400 documents, or 100 words. If 200 candidates remain and a split
removes 100, it saves 100 complete scores. It also adds 100 split-word visits and branch
work. For $\Delta F$ avoided scores, the local comparison is
\[
 \Delta F\,c_s>
 Wc_b+c_l\Delta V_l+\Delta T_{\rm branches}
       +\Delta T_{\rm row\ starts}+\Delta T_{\rm top}.
\]
The right side includes the changes caused by the extra split. Those changes need not all
be positive. The complete search must also retain the required recall; scoring fewer rows
by ending too early is insufficient.

If each split halves one candidate set, after $j$ cuts it contains approximately $n2^{-j}$
rows. Ignoring changes in queue and row-start costs gives the guide
\[
 T_B(j)\approx T_0+jWc_b+Wc_l+n2^{-j}c_s,\qquad
 j^*=\log_2\!\left(\frac{n\ln(2)c_s}{Wc_b}\right).
\]
This is an illustrative cost balance, not a recall estimate for these embeddings. If two
bitplanes are identical, splitting on the second after the first removes no additional rows.
Alternative branches also add paths. The experiment measures these effects rather than
assuming repeated halving.

B uses every cluster and probe option above, then varies local search: leaf size 128 with
node budgets 256 and 1024; leaf size 512 with unlimited work; and leaf size $N$, which permits
scoring an opened cluster without splitting. A node budget counts visited candidate sets,
including leaves. Leaf sizes 32 and 128 are also compared at budget 4096. Actual splits, final-mask reads and scored rows explain the resulting time.

The probe cutoff computed for A is a necessary routing condition for B using the same
layout. It may not suffice when local limits lose more results. Returned IDs determine
whether each complete setting qualifies.

\newpage
\subsection{Backward Walk}
The key must be wide enough to rule out some patterns. For the first saved query, the
ideal full score is 12.361. Even if all eight key bits mismatch, the remaining full-score
upper bound is 11.373. The 100th reference score is only 9.994. Thus even the least promising
eight-bit key cannot be discarded using this bound. All 200 queries have this problem at
key widths 1, 4, and 8.

At the same routing settings, changing key width from $h$ to $h'$ helps when
\[
 T_{\rm score}(F_h)-T_{\rm score}(F_{h'})>
 c_h(H_{h'}-H_h)+c_p(H_{+,h'}-H_{+,h})
       +T_{{\rm keys},h'}-T_{{\rm keys},h}.
\]
A smaller first range does not determine the answer. Search may need many additional keys,
including empty ones, before finding enough high-scoring documents or proving that remaining
keys cannot improve the result.

The experiment tries keys of 1, 4, 8, 16, 20, and 24 bits. For all 200 queries, the short
1-, 4-, and 8-bit keys cannot be discarded using the score bound, so their complete key spaces
are searched. Longer keys are compared with lookup limits of 4096 and 65,536.
The full 256-bit score remains unchanged. Directory memory is measured for
each index, and returned IDs establish recall. No $N/2^h$ occupancy assumption is used to
assign recall or declare a winning key width.

\subsection{Comparison}
Let $\Theta_a$ contain method $a$'s tested settings. Its final choice is
\[
 \theta_a^*=\underset{\theta\in\Theta_a}{\operatorname{argmin}}\;
 T_{a,\rm measured}(\theta)
 \quad\text{subject to}\quad
 \overline R_a(\theta)\ge\rho,\qquad
 M_{a,\rm process}(\theta)\le32\times10^9\text{ bytes}.
\]
\begin{center}\small
\begin{tabularx}{\linewidth}{@{}lYY@{}}\toprule
Method & Adjustable settings & Work that may be reduced\\\midrule
Original method & Cluster count $C$ and probes $P$ & Documents reaching the scorer\\
Bitplanes & $C$, $P$, leaf size and node budget & Scores, at the cost of mask and branch work\\
Backward Walk & $C$, $P$, key width $h$ and lookup limit & Scores, at the cost of lookups and range starts\\\bottomrule
\end{tabularx}\normalsize
\end{center}
The methods are compared after each has selected its own qualifying settings. Results below
the recall target remain visible, but cannot establish a speedup at that target. This finds
the best measured choices within the tested ranges on the unchanged query batch. It does not
prove that every possible index configuration has been exhausted.
